\chapter{Проекторно-тетраэдрическое отражение и запрет массовой поляризации}

Минимальный перенос оставил один непрерывный массовый модуль. Единственным
непроверенным внутренним кандидатом на его фиксацию были ранг-один проектор
и уже выведенная тетраэдрическая голономия $2\pi/3$. Настоящий гейт
проверяет, задают ли они не только ориентированную фазу, но и единственное
отражение, способное смешивать встречные хиральные направления.

\section{Антикруговая сверка}

Проект ранее доказал, что для каждого из четырёх проекторов $P_a$ и двух
ориентаций $\nu=\pm1$ существует канонический антисимметричный генератор
$\Omega(h_a)$ с точной голономией
\begin{equation}
 C_{a,\nu}=\exp\!\left(-\nu\frac{2\pi}{3}\Omega(h_a)\right).
 \label{eq:v5-tetra-holonomy}
\end{equation}
Все восемь ветвей совпадают с трёхциклами $S_4$ в машинной точности.

Но квантовое блуждание требует различать два типа данных. Ориентированная
комплексная структура задаёт фазовое вращение. Массовый член должен
обменивать две хиральные линии, то есть задавать отражение или
самосопряжённую инволюцию, антикоммутирующую с оператором ориентации.
Литература по хиральным квантовым блужданиям также рассматривает оператор
хиральной симметрии как дополнительную структуру эволюции; см.
\path{arXiv:1303.1199} и \path{arXiv:1502.02592}.

\section{Каноническая поперечная плоскость}

Пусть
\begin{equation}
 u=\frac12(1,1,1,1)^T,
 \qquad
 h_a=\frac2{\sqrt3}\left(e_a-\frac14\mathbf1\right).
\end{equation}
Проектор и аффинный синглет однозначно выделяют плоскость
\begin{equation}
 T_a=\{u,h_a\}^{\perp},
 \qquad
 P_{T_a}=I-uu^T-h_ah_a^T.
\end{equation}
На ней ориентационный генератор удовлетворяет
\begin{equation}
 \Omega_a^T=-\Omega_a,
 \qquad
 \Omega_a^2=-P_{T_a}.
 \label{eq:v5-transverse-complex-structure}
\end{equation}
Следовательно, ограничение $J_a=\Omega_a|_{T_a}$ является канонической
комплексной структурой двумерной вещественной плоскости.

Машинный аудит для всех четырёх проекторов даёт максимальный дефект
\eqref{eq:v5-transverse-complex-structure} меньше
$5.5\cdot10^{-16}$. Ограничение голономии равно
\begin{equation}
 C_{a,\nu}|_{T_a}
 =\cos\frac{2\pi}{3}\,I
 -\nu\sin\frac{2\pi}{3}\,J_a,
\end{equation}
а его собственные значения равны
\begin{equation}
 e^{2\pi i/3},\qquad e^{-2\pi i/3}.
\end{equation}
Тем самым плоскость, ориентация и фазовый угол действительно выведены.

\section{Почему проекторная инволюция не является массовым отражением}

Из кривизны прежней главы естественно возникает инволюция
\begin{equation}
 \mathcal J_a=2P_a-I_4.
\end{equation}
Она является единственным очевидным отражением, построенным только из
проектора. Однако на $T_a$ выполнено
\begin{equation}
 \mathcal J_a|_{T_a}=-I_{T_a}.
 \label{eq:v5-projector-involution-transverse}
\end{equation}
Поэтому она не выделяет прямую в поперечной плоскости и не меняет знак
$J_a$:
\begin{equation}
 [\mathcal J_a,J_a]=0.
\end{equation}
Она даёт общий знак или фазу, но не смешивает две сопряжённые хиральные
линии.

\section{Классификация допустимых отражений}

Выберем временный ортонормированный базис в $T_a$, в котором
\begin{equation}
 J=\begin{pmatrix}0&-1\\1&0\end{pmatrix}.
\end{equation}
Любая вещественная симметричная инволюция, меняющая ориентацию,
имеет вид
\begin{equation}
 R(\alpha)=
 \begin{pmatrix}
  \cos2\alpha&\sin2\alpha\\
  \sin2\alpha&-\cos2\alpha
 \end{pmatrix},
 \qquad \alpha\in\mathbb R/\pi\mathbb Z.
 \label{eq:v5-reflection-circle}
\end{equation}
Она удовлетворяет
\begin{equation}
 R(\alpha)^2=I,qquad
 R(\alpha)JR(\alpha)=-J,qquad
 \det R(\alpha)=-1.
\end{equation}
Таким образом, одна ориентированная плоскость несёт не одно отражение, а
целую окружность поляризаций.

Тетраэдрический стабилизатор выбранной вершины уменьшает эту окружность до
трёх естественных осей
\begin{equation}
 \alpha=0,\quad\frac\pi3,\quad\frac{2\pi}3.
\end{equation}
Но трёхцикл циклически переставляет их:
\begin{equation}
 R_0\longmapsto R_2\longmapsto R_1\longmapsto R_0.
\end{equation}
Ни одна ось не инвариантна, а их среднее равно нулю.

Более сильная линейная проверка даёт
\begin{equation}
 \{X:[X,C_3]=0,\ XJ+JX=0\}=\{0\}.
 \label{eq:v5-no-invariant-reflection}
\end{equation}
Матрица четырёх линейных условий на четыре компоненты $X$ имеет полный
ранг. Следовательно, ненулевого отражения, одновременно совместимого с
ориентацией и инвариантного относительно трёхцикла, не существует.

\begin{theorem}[Отсутствие канонической массовой поляризации]
Ранг-один проектор и тетраэдрическая голономия канонически задают
поперечную плоскость $T_a$, комплексную структуру $J_a$ и ориентацию
фазового переноса. Алгебра порождённых ими операторов на $T_a$ равна
\begin{equation}
 \operatorname{span}_{\mathbb R}\{I,J_a\}.
\end{equation}
Все её элементы коммутируют с $J_a$. Она не содержит отражения,
обменивающего две хиральные линии. Выбор такого отражения является новым
выбором поляризации.
\end{theorem}

\section{Заманчивый, но условный массовый масштаб}

Плоская связь прежней главы имеет плотность
\begin{equation}
 \mu_{\rm hol}=\frac{2\pi}{3L}.
\end{equation}
Если \emph{дополнительно} выбрать одно отражение $R(\alpha)$ и объявить
его генератором монеты на шаге длины $a$, то получится
\begin{equation}
 n=\cos\frac{2\pi a}{3L},
 \qquad
 m=\sin\frac{2\pi a}{3L},
\end{equation}
и формальный непрерывный массовый масштаб
\begin{equation}
 \lim_{a\to0}\frac{m}{a}=\frac{2\pi}{3L}.
 \label{eq:v5-conditional-holonomy-mass}
\end{equation}
Это важная подсказка: связь уже содержит размерную фазовую плотность.
Однако формула \eqref{eq:v5-conditional-holonomy-mass} не является выводом
массы, поскольку зависит от невыведенного $R(\alpha)$ и от отождествления
граничного параллельного переноса с монетой пространственного шага.

\begin{nogocriterion}[Запрет скрытой поляризации]
Нельзя использовать $2\pi/(3L)$ как массу, пока не получено единственное
функториальное отображение от проекторно-голономных данных к отражению
$R$ на хиральном дублете. Выбор одной из трёх осей стабилизатора или
произвольного $R(\alpha)$ является дополнительным секторным данным.
\end{nogocriterion}

\section{Корректная роль тетраэдрической голономии}

Отрицательный результат не уничтожает динамическое прозрение. Он уточняет
физический тип уже выведенного объекта. Голономия естественно задаёт
ориентированный фазовый перенос, а не массовое отражение. Поэтому следует
сохранить канонический безмассовый оператор и поместить тетраэдрический
трёхцикл в граничное условие.

В стандартном триплете трёхцикл имеет собственные значения
\begin{equation}
 1,\qquad e^{2\pi i/3},\qquad e^{-2\pi i/3}.
\end{equation}
Это немедленно предлагает следующий параметрически чистый тест: вычислить
спектр безмассового переноса на окружности с таким граничным условием.
Инвариантная линия может сохранить одну нулевую моду, тогда как две
сопряжённые ветви получают фиксированные сдвиги импульса на $\pm1/3$.

Это ещё не доказательство нейтринной природы моды. Но, в отличие от
массовой интерпретации, здесь голономия используется по своему
математическому назначению и не требует отражения или коэффициента Юкавы.

\begin{protocol}
\begin{enumerate}
 \item поперечная плоскость из $P_a$ и $u$:
 \textbf{получена};
 \item каноническая комплексная структура $J_a$:
 \textbf{получена};
 \item фазовая голономия $2\pi/3$:
 \textbf{получена};
 \item проекторная инволюция на $T_a$:
 \textbf{равна $-I$, массового смешивания не даёт};
 \item пространство отражений:
 \textbf{окружность};
 \item тетраэдрически естественные отражения:
 \textbf{три, циклически переставляются};
 \item инвариантное ненулевое отражение:
 \textbf{не существует};
 \item единственная массовая поляризация:
 \textbf{не выведена};
 \item минимальный путь вывода массы из семейной голономии:
 \textbf{закрыт};
 \item безмассовый голономный перенос:
 \textbf{сохранён};
 \item физическое замыкание:
 \textbf{не достигнуто}.
\end{enumerate}
\end{protocol}

Машинный сертификат сохранён в
\path{s2t_v5_rank_one_tetrahedral_transfer_reflection_gate.py} и
\path{s2t_v5_rank_one_tetrahedral_transfer_reflection_gate_results.json}.